\documentclass[french,a4paper]{article}

\usepackage[utf8]{inputenc}
\usepackage[T1]{fontenc}
\usepackage{lmodern}
\usepackage{geometry}
\usepackage{graphicx}
\usepackage{babel}
\usepackage{amsmath}
\usepackage{amsthm}
\usepackage{amssymb}
\usepackage{hyperref}
\usepackage{stmaryrd}
\usepackage{enumitem}
\usepackage{tcolorbox}
\usepackage{dsfont}
\usepackage{multirow}
\usepackage{etoc}
\usepackage{titlesec}
\usepackage{esint}
\usepackage{cancel}
\usepackage{mathdots}

\setlist[itemize]{label=$\bullet$}


\renewcommand{\P}{\mathbb{P}}
\newcommand{\N}{\mathbb{N}}
\newcommand{\Z}{\mathbb{Z}}
\newcommand{\Q}{\mathbb{Q}}
\newcommand{\R}{\mathbb{R}}
\newcommand{\C}{\mathbb{C}}
\newcommand{\K}{\mathbb{K}}
\renewcommand{\L}{\mathbb{L}}
\newcommand{\U}{\mathbb{U}}
\newcommand{\st}{^*}
\newcommand{\tq}{\middle|}
\newcommand{\inv}{^{-1}}
\newcommand{\E}{\mathbb{E}}
\newcommand{\F}{\mathbb{F}}
\newcommand{\V}{\mathbb{V}}
\renewcommand{\ker}{\text{Ker}}
\newcommand{\im}{\text{Im}}
\newcommand{\card}{\text{Card}}
\newcommand{\Tr}{\text{Tr}}

\theoremstyle{plain}
\newtheorem{thm}{Théorème}
\newtheorem*{ind}{Indication}

\theoremstyle{definition}
\newtheorem{dfn}{Definition}

\theoremstyle{remark}
\newtheorem*{rmq}{Remarque}
\newtheorem*{app}{Application}

\newtcolorbox[auto counter]{ex}[2][]{colback=white, colframe=green!50!white, coltitle=black, fonttitle=\bfseries, title={Exercice~\thetcbcounter~: #2}, after title={\hfill #1}}

\title{TD Continuité}


\begin{document}

\maketitle

\section{Fonctions d'une variable réelle}

\begin{ex}[Moulin.Cont.01]{$\circ$}
	Soit $f:\R_+\to\R$ continue telle que $\lvert f\rvert$ tende vers $+\infty$ en $+\infty$. Montrer que $f$ tend vers $+\infty$ en $+\infty$ ou $f$ tend vers $\infty$ en $+\infty$. Et sans l'hypothèse de continuité ?
\end{ex}
\begin{proof}
	Utiliser la continuité et faire un dessin. Sans continuité, on peut construire un contre exemple facilement.
\end{proof}

\begin{ex}[Moulin.Cont.02]{$\star$}
	Soit $f$ et $g$ deux fonctions continues de $[0,1]$ dans $[0,1]$ telles que $f\circ g=g\circ f$. Montrer que les graphes de $f$ et $g$ se coupent.
\end{ex}
\begin{proof}
	On raisonne par l'absurde et, quitte à échanger $f$ et $g$, d'après le théorème des bornes atteintes, il existe $m>0$ tel que $$\forall x\in[0,1],\ g(x)\geqslant f(x)+m$$ Par récurrence et grâce à la commutation, on montre que $$\forall n\in\N,\ \forall x\in[0,1],\ g^n(x)\geqslant f^n(x)+nm$$ (au sens de la composition itérée).  Par suite : $$\forall n\in\N, \ g^n(1)\geqslant nm$$ C'est absurde car le premier terme appartient toujours à $[0,1]$ et le deuxième diverge vers $+\infty$.
\end{proof}

\begin{ex}[Moulin.Cont.03]{$\Delta$}
	Soit $f:\R_+\to\R$ continue surjective. Montrer que $f\inv\left(\{a\}\right)$ est infini pour tout $a\in\R$.
\end{ex}
\begin{proof}
	Quitte à remplacer $f$ par $f-a$, ce qui ne change par le résultat, on peut suppose $a=0$. En raisonnant par l'absurde, si l'ensemble des zéros de $f$ était fini, $f$ serait de signe constant à partir d'un certain seuil, et ne pourrait donc pas être surjective car elle serait bornée avant d'après le théorème des bornes atteintes. C'est absurde.
\end{proof}

\begin{ex}[Moulin.Cont.04]{}
	Soit $f:\R_+\to\R$ continue et $l\in\overline{\R}$.
	\begin{enumerate}
		\item Montrer que si $f(x+1)-f(x)\xrightarrow[x\to+\infty]{}l$, alors $\displaystyle\frac{f(x)}{x}\xrightarrow[x\to+\infty]{}l$.
		\item Donner un contre-exemple sans l'hypothèse de continuité.
	\end{enumerate}
\end{ex}

\begin{ex}[Moulin.Cont.05]{}
	Soit $f:[0,1]\to\R$ telle que $f(0)<0<f(1)$. on suppose qu'il existe $g\in\mathcal{C}^0([0,1],\R)$ telle que $f+g$ soit décroissante. montrer qu'il existe $x\in[0,1]$ tel que $f(x)=0$.
\end{ex}

\begin{ex}[Moulin.Cont.06]{}
	Soit $f\in\mathcal{C}^0(\R,\R)$. Étudier les implications entre les assertions suivantes :
	\begin{enumerate}[label=(\roman*)]
		\item $f$ est uniformément continue ;
		\item $f$ admet des limites finies en $\pm\infty$ ;
		\item $f(x)=O(x)$ quand $x\to\pm\infty$ ;
		\item il existe $(A,B)\in\R_+^2$ tel que $\forall x\in\R, \lvert f(x)\rvert\leqslant A\lvert x\rvert+B$.
	\end{enumerate}
\end{ex}

\begin{ex}[Moulin.Cont.07]{}
	Soit $X$ une partie de $A$ dense dans $A$ et $f:X\to F$ continue en tout point de $X$. On suppose que $f$ admet une limite finie en tout point de $A\setminus X$. Montrer que $f$ admet un prolongement $\tilde{f}:A\to F$ continu en tout point de $A$.
\end{ex}

\section{Étude locale}

\begin{ex}[Moulin.Cont.08]{$\circ$}
	\begin{enumerate}
		\item Montrer que toute fonction est continue en un point isolé de son ensemble de définition $A$, c'est-à-dire en un point $a\in A$ non adhérent à $A\setminus\{a\}$. En particulier, une suite à valeurs dans un espace vectoriel normé $E$ est une application continue de $\N$ dans $E$.
		\item Que dire d'une application continue de $\R$ dans $\N$ ?
	\end{enumerate}
\end{ex}

\begin{ex}[Moulin.Cont.09]{$\circ$}
	Soit $f:A\to\displaystyle\prod_{i=1}^{n}E_i$ où les $(E_i,N_i)$ sont des espaces vectoriels normés. Soit $a$ un point adhérent à $A$. Montrer que s'il existe $i$ tel que $\underset{x\to a}{\lim} f_i(x)=\infty$, alors $\underset{x\to a}{\lim} f(x)=\infty$. Réciproque ?
\end{ex}

\begin{ex}[Moulin.Cont.10]{$\star$}
	Soit $f$ une fonction continue de $E=\R^2\setminus{(0,0)}$ dans $\R$ et $a\in\overline{\R}$. A-t-on :
	\begin{enumerate}
		\item $\left(\forall u\in E,\ \underset{t\to+\infty}{\lim}f(tu)=a\right)\implies \underset{\infty}{\lim} f=a$ ?
		\item $\left(\forall u\in E,\ \underset{t\to0}{\lim}f(tu)=a\right)\implies \underset{0}{\lim} f=a$ ?
	\end{enumerate}
\end{ex}

\section{Continuité globale}

\begin{ex}[Moulin.Cont.11]{}
	L'application qui à $(x,y)$ associe $x(1-y)$ si $x\geqslant y$ et $y(1-x)$ sinon est-elle continue sur $\R^2$ ?
\end{ex}

\begin{ex}[Moulin.Cont.12]{$\star$}
	Soit $f:\R^2\to\R$. On suppose $f(a,\cdot)$ et $f(\cdot,b)$ continues et monotones pour tout $(a,b)\in\R^2$. Montrer que $f$ est continue en tout point de $\R^2$.
\end{ex}
\begin{proof}
	Faire un grand dessin. On prend $r_0$ pour la continuité en $(a,b)$, $r_1$ pour la continuité de $f(\cdot,b)$ et $r_2$ pour celle de $f(a,\cdot)$ et le minimum de $r_0$, $r_1$ et $r_2$ convient par monotonies.
\end{proof}

\begin{ex}[Moulin.Cont.13]{}
	Montrer que la fonction qui à $(x,y)\in(\R_+)^2\setminus\{(0,0)\}$ associe $x^y$ est continue. Admet-elle une limite en $(0,0)$ ? Admet-elle une limite à l'infini ?
\end{ex}

\begin{ex}[Moulin.Cont.14]{$\Delta$}
	Soit $n\in\N\st$. Montrer que les fonctions : $$m:(x_1,\ldots,x_n)\mapsto\min(x_1,\ldots,x_n)\ \ \text{et}\ \ M:(x_1,\ldots,x_n)\mapsto\max(x_1,\ldots,x_n)$$ sont continues.
\end{ex}
\begin{proof}
	On montre aisément que le minimum et le maximum sont associatifs au sens suivant : $$\begin{cases}
		\min(x_0,x_1,\ldots,x_n)=\min\left[x_0,\min(x_1,\ldots,x_n)\right] \\
		\max(x_0,x_1,\ldots,x_n)=\max\left[x_0,\max(x_1,\ldots,x_n)\right]
	\end{cases}$$ Par une récurrence immédiate, il suffit simplement de montrer la continuité des application : $$\min:(x,y)\mapsto\min(x,y) \ \ \text{et} \ \ \max:(x,y)\mapsto\max(x,y)$$ Or, on sait que $$\min(x,y)=\frac{x+y}{2}-\frac{\lvert y-x\rvert}{2} \ \ \text{et} \ \ \max(x,y)=\frac{x+y}{2}+\frac{\lvert y-x\rvert}{2}$$ On déduit de ces deux expressions que ces fonctions sont continues (caractère polynomial et composition avec la valeur absolue), ce qui conclut.
\end{proof}

\begin{ex}[Moulin.Cont.15]{}
	Étudier le prolongement par continuité de la fonction $f$ définie sur $\R^2\setminus\{(0,0)\}$ par $$f(x,y)=\frac{e^y-e^x}{y-x}$$
\end{ex}

\begin{ex}[Moulin.Cont.16]{$\Delta$}
	Soit $f:\R\mapsto\C$. On note $\Delta=\{(x,x)\mid x\in\R\}$. Montrer que la fonction : $$\begin{array}{lrcl}
		F:&\R^2\setminus\Delta&\longrightarrow&\C\\
		&(x,y)&\longmapsto&\displaystyle\frac{f(x)-f(y)}{x-y}
	\end{array}$$ est prolongeable en une fonction continue sur $\R^2$ si, et seulement si, $f$ est de classe $\mathcal{C}^1$.
\end{ex}
\begin{proof}
	Dans le sens direct, un tel prolongement est nécessairement un prolongement par la dérivée (en prenant une continuité partielle). puis la continuité selon $\Delta$ fournit la continuité de la dérivée. Dans le sens réciproque, considérer le seul prolongement possible, \textit{ie} celui par la dérivée sur $\Delta$. Utiliser ensuite l'\textbf{inégalité des accroissements finis} (la TAF est faux dans $\C$) pour prouver que ce prolongement est bien continu selon $\R^2\setminus\Delta$ (et il est continu selon $\Delta$ car $f$ est de classe $\mathcal{C}^1$). On pourra utiliser le fait que $$F(x,y)=\int_{0}^{1}f'(x+t(y-x))\ dt$$ Cela permet de condenser les deux informations en une seule fonction (cette astuce est utile dans le chapitre Intégrales à paramètres).
\end{proof}

\begin{ex}[Moulin.Cont.17]{}
	Soit $\varphi:\R\to\R$. On considère $f$ de $\R^2$ dans $\R$ qui à $(x,y)$ associe $\displaystyle\frac{1-\cos(\sqrt{\lvert xy\rvert})}{\lvert y\rvert}$ si $y\ne 0$ et $\varphi(x)$ sinon. Déterminer une condition nécessaire et suffisante portant sur $\varphi$ pour que $f$ soit continue.
\end{ex}

\begin{ex}{$\Delta$}
	Montrer que toute application continue sur $\R$ qui admet des limites finies en $-\infty$ et en $+\infty$ est uniformément continue.
\end{ex}
\begin{proof}
	Faire un dessin pour essayer de comprendre ce qui suit ! Prendre $\varepsilon>0$ et appliquer la définition des limites avec $\varepsilon'=\displaystyle\frac{\varepsilon}{2}$ pour obtenir $A$ (pour $-\infty$) et $B$ (pour $+\infty$). Sans perte de généralité, on peut supposer $B\geqslant A+2$. Ensuite, on applique le théorème de Heine à la restriction de la fonction au segment $[A-1,B+1]$ et sans perte de généralité, quitte à diminuer $\delta$, on suppose $0<\delta\leqslant1/2$. On vérifie alors qu'un tel $\delta$ convient : avant $A-1$, pas de problème par définition de $A$ et avec l'inégalité triangulaire, idem après $B+1$ ; ensuite, au niveau de $A-1$ et $B+1$, OK car on a supposé $\delta\leqslant1/2$ ; entre $A-1$ et $B+1$, OK par le théorème de Heine. Et il n'y a pas d'autre cas puisqu'on a imposé $B\geqslant A+2$.
\end{proof}

\begin{ex}[Moulin.Cont.18]{$\Delta$}
	Montrer que toute application continue périodique sur $\R$ est uniformément continue.
\end{ex}
\begin{proof}
	Appliquer le théorème de Heine sur $[0,2T]$ puis par périodicité on se ramène $x$ et $y$ dans $[0,2T]$ en ajoutant des multiples de $T$.
\end{proof}

\begin{ex}[Moulin.Cont.19]{}
	Étudier la continuité de l'application $f\mapsto\varphi\circ f$ de $\left(\mathcal{C}^0([a,b],\R),\mathcal{N}_\infty\right)$ dans lui même, où $\varphi$ est continue sur $\R$. A quelle condition sur $\varphi$ cette application est-elle uniformément continue ? Lipschitzienne ?
\end{ex}

\begin{ex}[Moulin.Cont.20]{}
	Montrer qu'une partie de $\R$ est fermée si, et seulement si, c'est l'ensemble des zéros d'une fonction continue de $\R$ dans $\R$.
\end{ex}

\begin{ex}[Moulin.Cont.21]{}
	Soit $E$ et $F$ deux espaces vectoriels normés et $A$ une partie de $E$.
	\begin{enumerate}
		\item Montrer que si une application $f:A\to F$ est continue, son graphe est fermé dans $A\times F$.
		\item La réciproque est-elle vraie en général ?
		\item $\star$ Montrer que la réciproque est vraie si $F=\R$ et si $f$ est bornée.
	\end{enumerate}
\end{ex}

\begin{ex}[Moulin.Cont.22]{$\star\star$}
	Montrer que l'on peut définir une application $\varphi$ de $\R^2$ dans $\R$ en associant à $(x,y)\in\R^2$ la plus grand racine réelle du polynôme $X^5+xX+y$. En quels points $\varphi$ est-elle continue ?
\end{ex}

\begin{ex}[Moulin.Cont.23]{$\star\star$ Continuité du passage à l'inverse}
	Soit $E$ la partie de $\left(\mathcal{C}([0,1],\R),\mathcal{N}_\infty\right)$ constitué des fonctions strictement croissantes vérifiant $f(0)=0$ et $f(1)=1$.
	\begin{enumerate}
		\item Montrer que l'application $\varphi:f\mapsto f\inv$ est continue sur $E$.
		\item $\varphi$ est-elle uniformément continue sur $E$ ?
	\end{enumerate}
\end{ex}
\begin{proof}
	Exercice très difficile. Merci Pierre pour les travaux.
	\begin{enumerate}
		\item On se donne $f$ et $g$ dans $E$. Alors : \begin{align*}
			\mathcal{N}_\infty(g\inv-f\inv)&=\underset{x\in[0,1]}{\sup}\lvert g\inv(x)-f\inv(x)\rvert \\
			&=\underset{t\in[0,1]}{\sup}\lvert t-f\inv(g(t))\rvert\\
			&=\mathcal{N}_\infty(\text{Id}-f\inv\circ g)
		\end{align*} Or, $f\inv$ est continue sur $[0,1]$ donc uniformément continue d'après le théorème de Heine. Soit $\varepsilon>0$, et prenons alors $\delta>0$ tel que $$\forall (x,y)\in[0,1]^2,\ \lvert x-y\rvert\leqslant\delta \implies\lvert f\inv(x)-f\inv(y)\rvert\leqslant\varepsilon$$ Soit $t\in[0,1]$. On pose $x=f(t)$ et $y=g(t)$. Alors, si $\mathcal{N}_\infty(g-f)\leqslant\delta$, on en déduit que $$\lvert t-f\inv(g(t))\rvert=\lvert f\inv(x)-f\inv(y)\rvert\leqslant\varepsilon$$ Par conséquent, $$\mathcal{N}_\infty(\text{Id}-f\inv\circ g)\leqslant\varepsilon$$ puisque $f$ est bijective.
		\item On va montrer que $f$ n'est pas uniformément continue. On considère $f$ et $g$ affine par morceaux vérifiant les conditions et telles que $f(1/2)=\alpha$ et $g(1/2)=\beta$. Alors : $$\mathcal{N}_\infty(f-g)=\lvert\alpha-\beta\rvert$$ Et on vérifie en calculant les inverses que $$\mathcal{N}_\infty(f\inv-g\inv)=\frac{1}{2}-\frac{\beta}{2\alpha}$$ On considère alors les suites $\alpha_n=\displaystyle\frac{1}{2^n}$ et $\beta_n=\displaystyle\frac{1}{2^{n+1}}$ ainsi que les fonctions définies comme $f$ et $g$ précédemment. Alors $\lvert\alpha_n-\beta_n\rvert\to0$ mais pourtant $$\frac{1}{2}-\frac{\beta_n}{2\alpha_n}=\frac{1}{4}\to\frac{1}{4}\ne0$$
		Donc $\varphi$ n'est pas uniformément continue par négation de la caractérisation séquentielle de l'uniforme continuité (déjà pas très au programme).
	\end{enumerate}
\end{proof}

\begin{ex}[Moulin.Cont.24]{Fonction en escalier sur des fermés}
	Soit $E$ un espace vectoriel normé.
	\begin{enumerate}
		\item Soit $F_0$ et $F_1$ deux fermés disjoints non vides. En utilisant les fonction $x\mapsto d(x,F_i)$, trouver une fonction continue égale à $1$ sur $F_1$ et nulle sur $F_0$.
		\item Soit $F_1,\ldots,F_n$ des fermés deux à deux disjoints. Montrer qu'il existe, pour tout $(a_1,\ldots,a_n)\in\R^n$, une fonction $f$ continue telle que $$\forall i\in\llbracket1,n\rrbracket,\ \forall x\in F_i,\ f(x)=a_i$$
	\end{enumerate}
\end{ex}
\begin{proof}
	Remarquer qu'on a une sorte d'interpolation de Lagrange à mener.
	\begin{enumerate}
		\item La fonction $$f:x\mapsto\frac{d(x,F_1)}{d(x,F_0)+d(x,F_1)}$$ convient. Si on pose $$U_0=f\inv\left(]-\infty ;1/4[\right) \ \ \text{et} \ \ U_2=f\inv\left(]3/4;+\infty[\right)$$ on vérifie sans peine que $F_0\subset U_0$ et $F_1\subset U_1$ et que $U_0$ et $U_1$ sont des ouverts comme images réciproques d'ouverts.
		\item On cherche à généraliser ce qui précède. La fonction $$f:x\mapsto\frac{\sum_{i=1}^{n}a_i\prod_{j\ne i}d(x,F_i)}{\sum_{i=1}^{n}\prod_{j\ne i}d(x,F_i)}$$ convient de même.
	\end{enumerate}
\end{proof}

\section{Applications linéaires continues}

\begin{ex}[Moulin.Cont.25]{$\Delta$ La dérivation n'est pas continue sur $\mathcal{C}^\infty$}
	Montrer qu'il n'existe pas de norme sur $\mathcal{C}^{\infty}([0,1],\R)$ pour laquelle la dérivation soit continue.
\end{ex}
\begin{proof}
	Raisonner par l'absurde et considérer alors une constante $C\geqslant 0$ telle que $$\forall f,\ \lVert D(f)\rVert\leqslant C\lVert f\rVert$$ Mais on connaît de beaux vecteurs propres de la dérivation : les fonctions exponentielles. Si on considère $$f:x\mapsto e^{(C+1)x}$$ alors $D(f)=(C+1)f$ et $f$ n'est pas la fonction nulle donc on a une absurdité.
\end{proof}
\begin{rmq}
	Lorsqu'on travaille avec la dérivation sur $\mathcal{C}^\infty$, bien se souvenir qu'on a énormément de vecteurs propres.
\end{rmq}

\begin{ex}[Moulin.Cont.26]{}
	\begin{enumerate}
		\item Soit $D$ l'opérateur de dérivation sur $\R[X]$. Donner un exemple de norme pour laquelle $D$ est continu, et un exemple de norme pour laquelle $D$ n'est pas continu.
		\item Soit $M$ l'endomorphisme de $E$ qui au polynôme $P$ associe le polynôme $XP$. Existe-t-il une norme sur $E$ qui rende simultanément $D$ et $M$ continus ?
	\end{enumerate}
\end{ex}

\begin{ex}[Moulin.Cont.27]{}
	Déterminer les valeurs de $\alpha\in\C$ l'application $P\mapsto P(\alpha)$ est-elle continue lorsque l'on munit $\C[X]$ des normes définies, pour $P=\displaystyle\sum_{n=0}^{+\infty}a_nX^n$, par : $$\lVert P\rVert_1=\sum_{n=0}^{+\infty}\lvert a_n\rvert \ \ \lVert P\rVert_\infty=\underset{n\in\N}{\sup}\lvert a_n\rvert\ \ \text{et}\ \ \lVert P\rVert_2=\left(\sum_{n=0}^{+\infty}\lvert a_n\rvert^2\right)^{1/2}$$
\end{ex}

\begin{ex}[Moulin.Cont.28]{}
	Soit $E=\mathcal{C}^0([0,1],\R)$. Pour $\rho\in E$, on définit la forme linéaire : $$\varphi: f\mapsto\int_{0}^{1}f(t)\rho(t)\ dt$$ L'application $\varphi$ est-elle continue lorsque l'on munit $E$ de chacune des trois normes $\mathcal{N}_1$, $\mathcal{N}_2$ et $\mathcal{N}_\infty$ ? Calculer alors sa norme subordonnée. Est-elle atteinte ? 
\end{ex}

\begin{ex}[Moulin.Cont.29]{}
	Montrer que $f\mapsto\displaystyle\int_{0}^{1/2}f(t)\ dt-\int_{1/2}^{1}f(t)\ dt$ définit une forme linéaire continue sur $\left(\mathcal{C}^0([0,1]),\mathcal{N}_\infty\right)$. Calculer alors sa norme subordonnée. Est-elle atteinte ?
\end{ex}

\begin{ex}[Moulin.Cont.30]{$\star$}
	On note $\mathcal{C}$ l'espace vectoriel des suites réelles convergentes et $\mathcal{C}_0$ le sous-espace vectoriel de $\mathcal{C}$ constitué des suites convergeant vers $0$. On les munit de la norme $\lVert\cdot\rVert_\infty$.
	\begin{enumerate}
		\item Construire un homéomorphisme linéaire de $\mathcal{C}$ sur $\mathcal{C}_0$.
		\item Montrer qu'il n'existe pas d'isométrie linéaire bijective de $\mathcal{C}$ sur $\mathcal{C}_0$. On pourra s'intéresser aux boules sur $\mathcal{C}$ de centre $e$ et $-e$ et de rayon $1$, pù $e$ est la suite constante $(1)_{n\in\N}$.		
	\end{enumerate}
\end{ex}
\begin{proof}
	Exercice difficile. Merci Tim pour la correction.
	\begin{enumerate}
		\item Si $(u_n)$ a pour limite $l$, on note $f((u_n))$ la suite $(v_n)$ telle que $v_0=l$ et $$\forall n\in\N\st,\ v_{n}=u_{n-1}-l$$ Alors $f$ est évidemment linéaire et sa bijection réciproque est donnée par la fonction qui à $(v_n)$ associe la suite $(u_n)=(v_n+v_0)$. Les deux applications linéaires sont continues car on montre aisément que leurs normes triples sont toutes les deux inférieures à $2$.
		\item Cette question n'a pas de rapport avec la précédente. Notons que l'intersection des deux boules fermées de l'énoncé est réduite à la suite nulle. On raisonne par l'absurde en supposant qu'il existe une telle isométrie $\varphi$ linéaire bijective de $\mathcal{C}$ sur $\mathcal{C}_0$. Alors, puisque c'est une isométrie : $$\varphi\left(B_f(e,1)\right)=B_f(\varphi(e),1)$$ et de même pour l'autre boule. On prend alors $n_0\in\N$ tel que $$\lvert\varphi(e)_{n_0}\rvert\leqslant\frac{1}{2}$$ On définit alors la suite $u$ par $u_n=0$ si $n\ne n_0$ et $u_{n_0}=\displaystyle\frac{1}{2}$. Alors : $$\lVert u-\varphi(e)\rVert\leqslant 1\ \ \text{et} \ \ \lVert u+\varphi(e)\rVert\leqslant 1$$ Donc $$\{(0)\}\subsetneq \varphi\left(B_f(-e,1)\right)\cap \varphi\left(B_f(e,1)\right)$$ ce qui est absurde.
	\end{enumerate}
\end{proof}

\begin{ex}[Moulin.Cont.31]{Supplémentaires topologiques}
	Soit $E$ un espace vectoriel normé et une décomposition $E=F\oplus G$.
	\begin{enumerate}
		\item Montrer l'équivalence des assertions suivantes :
		\begin{enumerate}[label=(\roman*)]
			\item l'application de $F\times G$ dans $E$ qui à $(x,y)$ associe $x+y$ est un homéomorphisme ;
			\item la projection sur $F$ parallèlement à $G$ est continue ;
			\item $\exists\alpha>0\ : \ \forall x\in G,\ \alpha\lVert x\rVert \leqslant d(x,F)$ ;
			\item $\exists\beta>0\ : \ \forall x\in F,\ \beta\lVert x\rVert \leqslant d(x,G)$ ;
			\item $x\mapsto d(x,F)+d(x,G)$ est une norme sur $E$ équivalente à $\lVert\cdot\rVert$.
		\end{enumerate}
		\item Lorsque ces conditions sont vérifiées, montrer que $F$ et $G$ sont fermés dans $E$.
		\item $\star$ Donner un exemple de décomposition $E=F\oplus G$ où $F$ et $G$ sont fermés dans $E$ et ne vérifient pas les conditions de la première question.
	\end{enumerate}
\end{ex}

\begin{ex}[Moulin.Cont.32]{}
	Soit $E=\mathcal{C}^0([0,1],\R)$. Une forme linéaire $L:E\to\R$ est dite positive lorsque $f\geqslant 0\implies L(f)\geqslant0$.
	\begin{enumerate}
		\item Montrer qu'une forme linéaire positive est continue lorsque l'on munit $E$ de la norme uniforme. Quelle est sa norme subordonnée ?
		\item Le résultat reste-t-il vrai si l'on remplace la norme uniforme par la norme $L^1$ ?
	\end{enumerate}
\end{ex}

\section{Normes subordonnées}

\begin{ex}[Moulin.Cont.33]{$\Delta$}
	Soit $f$ et $g$ deux endomorphismes d'un espace vectoriel normé $E\ne\{0\}$ tels que $f\circ g-g\circ f=\text{Id}_E$.
	\begin{enumerate}
		\item Calculer $f\circ g^n-g^n\circ f$ pour tout $n\in\N$ et en déduire que $f$ et $g$ ne peuvent pas être simultanément continues.
		\item Donner deux bonnes raisons de la non-existence d'un tel couple $(f,g)$ lorsque $E$ est de dimension finie.
	\end{enumerate}
\end{ex}
\begin{proof}
	A mettre en regard du DM sur l'étude de ce crochet de Lie.
	\begin{enumerate}
		\item On remarque que $$f\circ g^{n+1}-g^{n+1}\circ f=(f\circ g^n-g^n\circ f)\circ g+g^n\circ \underbrace{(f\circ g-g\circ f)}_{=\text{Id}_E}$$ ce qui permet de montrer par une récurrence immédiate que $$\forall n\in\N,\ f\circ g^n-g^n\circ f=ng^{n-1}$$ En raisonnant par l'absurde et en supposant  $f$ et $g$ continues, si on note $N$ la norme subordonnée, on obtient :  $$nN(g^{n-1})\leqslant2N(f)N(g)N(g^{n-1})$$ Donc $N(g^{n-1})=0$ APCR sinon $2N(f)N(g)$ exploserait. Et si on suppose $g^{n}=0$ pour $n\geqslant 1$, alors la première relation obtenue nous donne $g^{n-1}=0$. Donc par récurrence $\text{Id}_E=0$, ce qui est absurde car $E\ne\{0\}$.
		\item En dimension finie :
		\begin{itemize}
			\item $f$ et $g$ sont linéaires donc continues, donc par contraposée de la première question, l'existence d'un tel couple est impossible ;
			\item on aurait $\Tr(\text{Id}_E)=\Tr(f\circ g-g\circ f)=0$ donc $E\ne\{0\}$ ce qui n'est pas le cas ici.
		\end{itemize}
	\end{enumerate}
\end{proof}

\begin{ex}[Moulin.Cont.34]{}
	Déterminer, sur $\mathcal{M}_{n,p}(\K)$ identifié canoniquement à $\mathcal{L}(\K^p,\K^n)$, les normes subordonnées au normes $\lVert\cdot\rVert_\infty$ de $\K^p$ et $\K^n$. Même question pour les normes $\lVert\cdot\rVert_1$. Même question si l'on munit $\K^p$ de la norme $\lVert\cdot\rVert_1$ et $\K^n$ de la norme $\lVert\cdot\rVert_\infty$.
\end{ex}

\begin{ex}[Moulin.Cont.35]{}
	Soit $n\in\N\st$, $\lVert\cdot\rVert_\infty$ la norme du maximum des modules des coefficients sur $\mathcal{R}^n$ et $N$ sa norme subordonnée sur $\mathcal{M}_n(\R)$. On fixe $A=(a_{i,j})_{1\leqslant i,j\leqslant n}\in\mathcal{M}_n(\R)$.
	\begin{enumerate}
		\item $\Delta$ Exprimer $N(A)$ au moyen des quantités $L_i(A)=\displaystyle\sum_{j=1}^{n}\lvert a_{i,j}\rvert$ pour $1\leqslant i\leqslant n$.
		\item On suppose que, pour tout $X$ de $\R^n$, on a : $\lVert AX\rVert_\infty=\lVert X\rVert_\infty$. Montrer que $A$ est le produit d'une matrice diagonale à termes diagonaux valant $1$ ou $-1$ et d'une matrice de permutation.
		\item On suppose $A$ inversible. On définit $\mu=\inf\{\lVert AX\rVert_\infty\ : \ X\in\R^n,\ \lVert X\rVert_\infty=1\}$. Montrer que $\mu>0$ et que $$\mu=\frac{1}{N(A\inv)}$$
		\item $\star\star$ Toujours si $A$ est inversible, montrer que $\mu$ est la distance de $A$ à l'ensemble des matrices non inversibles de $\mathcal{M}_n(\R)$.
	\end{enumerate}
\end{ex}

\begin{ex}[Moulin.Cont.36]{}
	Soit $(E,\lVert\cdot\rVert)$ une espace vectoriel normé et $N$ la norme subordonnée sur $\mathcal{L}_c(E)$.
	\begin{enumerate}
		\item Soit $p$ et $q$ deux projecteurs continus avec $0\leqslant\text{rg}(q)<\text{rg}(p)\leqslant+\infty$. Montrer que $N(p-q)\geqslant 1$.
		\item Soit $P:[0,1]\to\mathcal{L}_c(E)$ continue. On suppose que, pour tout $t\in[0,1]$, $P(t)$ est un projecteur continu de rang fini. Montrer que l'application $t\mapsto\text{rg}(P(t))$ est constante.
	\end{enumerate}
\end{ex}

\begin{ex}[Moulin.Cont.37]{}
	Soit $n\in\N\st$. On munit $\R_{n-1}[X]$ de la norme $N(P)=\displaystyle\sum_{k=1}^{n}\lvert P(k)\rvert$. Calculer la norme triple de la forme linéaire $P\mapsto P(0)$.
\end{ex}

\begin{ex}[Moulin.Cont.38]{$\star$ Théorème de Hahn-Banach en dimension finie}
	Soit $E$ un $\R$-espace vectoriel normé de dimension finie, $F$ un sous-espace vectoriel de strict de $E$ et $\varphi$ une forme linéaire sur $F$ dont la norme triple est égale à $1$.
	\begin{enumerate}
		\item Soit $a\in E$? Montrer : $$\forall (x,y)\in F^2,\ \varphi(x)-\lVert x-a\rVert\leqslant\varphi(y)+\lVert a-y\rVert$$
		\item Montrer qu'il existe un prolongement de $\varphi$ en une forme linéaire $\tilde{\varphi}$ sur $E$ dont la norme triple est égale à $1$.
	\end{enumerate}
\end{ex}

\end{document}